\documentclass[10pt,a4paper]{amsart}
\usepackage{amsmath,amssymb,amsthm,graphicx}
\usepackage{verbatim}
\newenvironment{reference}{\comment}{\endcomment}
%\newenvironment{reference}{}{}
\newenvironment{slogan}{\comment}{\endcomment}
\newenvironment{history}{\comment}{\endcomment}

% For commutative diagrams we use Xy-pic
\usepackage[all]{xy}

% We use 2cell for 2-commutative diagrams.
\xyoption{2cell}
% Theorem environments.
%
\theoremstyle{plain}
\newtheorem{theorem}[subsection]{Theorem}
\newtheorem{proposition}[subsection]{Proposition}
\newtheorem{lemma}[subsection]{Lemma}

\theoremstyle{definition}
\newtheorem{definition}[subsection]{Definition}
\newtheorem{example}[subsection]{Example}
\newtheorem{exercise}[subsection]{Exercise}
\newtheorem{situation}[subsection]{Situation}

\theoremstyle{remark}
\newtheorem{remark}[subsection]{Remark}
\newtheorem{remarks}[subsection]{Remarks}

\numberwithin{equation}{subsection}

% Macros
%
\def\lim{\mathop{\mathrm{lim}}\nolimits}
\def\colim{\mathop{\mathrm{colim}}\nolimits}
\def\Spec{\mathop{\mathrm{Spec}}}
\def\Hom{\mathop{\mathrm{Hom}}\nolimits}
\def\Ext{\mathop{\mathrm{Ext}}\nolimits}
\def\SheafHom{\mathop{\mathcal{H}\!\mathit{om}}\nolimits}
\def\SheafExt{\mathop{\mathcal{E}\!\mathit{xt}}\nolimits}
\def\Sch{\mathit{Sch}}
\def\Mor{\mathop{\mathrm{Mor}}\nolimits}
\def\Ob{\mathop{\mathrm{Ob}}\nolimits}
\def\Sh{\mathop{\mathit{Sh}}\nolimits}
\def\NL{\mathop{N\!L}\nolimits}
\def\CH{\mathop{\mathrm{CH}}\nolimits}
\def\proetale{{pro\text{-}\acute{e}tale}}
\def\etale{{\acute{e}tale}}
\def\QCoh{\mathit{QCoh}}
\def\Ker{\mathop{\mathrm{Ker}}}
\def\Im{\mathop{\mathrm{Im}}}
\def\Coker{\mathop{\mathrm{Coker}}}
\def\Coim{\mathop{\mathrm{Coim}}}

% Boxtimes
%
\DeclareMathSymbol{\boxtimes}{\mathbin}{AMSa}{"02}

%
% Macros for moduli stacks/spaces
%
\def\QCohstack{\mathcal{QC}\!\mathit{oh}}
\def\Cohstack{\mathcal{C}\!\mathit{oh}}
\def\Spacesstack{\mathcal{S}\!\mathit{paces}}
\def\Quotfunctor{\mathrm{Quot}}
\def\Hilbfunctor{\mathrm{Hilb}}
\def\Curvesstack{\mathcal{C}\!\mathit{urves}}
\def\Polarizedstack{\mathcal{P}\!\mathit{olarized}}
\def\Complexesstack{\mathcal{C}\!\mathit{omplexes}}
% \Pic is the operator that assigns to X its picard group, usage \Pic(X)
% \Picardstack_{X/B} denotes the Picard stack of X over B
% \Picardfunctor_{X/B} denotes the Picard functor of X over B
\def\Pic{\mathop{\mathrm{Pic}}\nolimits}
\def\Picardstack{\mathcal{P}\!\mathit{ic}}
\def\Picardfunctor{\mathrm{Pic}}
\def\Deformationcategory{\mathcal{D}\!\mathit{ef}}
\usepackage{cleveref}
\crefname{originallemma}{original lemma}{original lemmas}
\Crefname{originallemma}{Original lemma}{Original lemmas}
\makeatletter
\newlabel{decent-spaces-lemma-dimension-local-ring-quasi-finite}{{0ABW}{source}{}{sourcepointer-1}{}}
\newlabel{decent-spaces-lemma-dimension-local-ring-quasi-finite@cref}{{[originallemma][0][]0ABW}{[1][0][]source}{}{}{}}
\makeatother
\crefname{originallemma}{original lemma}{original lemmas}
\Crefname{originallemma}{Original lemma}{Original lemmas}
\makeatletter
\newlabel{more-groupoids-lemma-find-affine-codimension-1}{{0ABI}{source}{}{sourcepointer-2}{}}
\newlabel{more-groupoids-lemma-find-affine-codimension-1@cref}{{[originallemma][0][]0ABI}{[1][0][]source}{}{}{}}
\makeatother
\crefname{originallemma}{original lemma}{original lemmas}
\Crefname{originallemma}{Original lemma}{Original lemmas}
\makeatletter
\newlabel{spaces-cohomology-lemma-image-affine-finite-morphism-affine-Noetherian}{{07VP}{source}{}{sourcepointer-3}{}}
\newlabel{spaces-cohomology-lemma-image-affine-finite-morphism-affine-Noetherian@cref}{{[originallemma][0][]07VP}{[1][0][]source}{}{}{}}
\makeatother
\crefname{originalproposition}{original proposition}{original propositions}
\Crefname{originalproposition}{Original proposition}{Original propositions}
\makeatletter
\newlabel{spaces-limits-proposition-there-is-a-scheme-finite-over}{{09YC}{source}{}{sourcepointer-4}{}}
\newlabel{spaces-limits-proposition-there-is-a-scheme-finite-over@cref}{{[originalproposition][0][]09YC}{[1][0][]source}{}{}{}}
\makeatother
\crefname{originaldefinition}{original definition}{original definitions}
\Crefname{originaldefinition}{Original definition}{Original definitions}
\makeatletter
\newlabel{spaces-properties-definition-dimension-local-ring}{{04NA}{source}{}{sourcepointer-5}{}}
\newlabel{spaces-properties-definition-dimension-local-ring@cref}{{[originaldefinition][0][]04NA}{[1][0][]source}{}{}{}}
\makeatother
\crefname{originallemma}{original lemma}{original lemmas}
\Crefname{originallemma}{Original lemma}{Original lemmas}
\makeatletter
\newlabel{topology-lemma-closed-morphism-quotient-topology}{{0AAU}{source}{}{sourcepointer-6}{}}
\newlabel{topology-lemma-closed-morphism-quotient-topology@cref}{{[originallemma][0][]0AAU}{[1][0][]source}{}{}{}}
\makeatother
\title{A separated locally Noetherian algebraic space is a scheme at points whose local ring has dimension at most one}
\author{The Stacks Project Authors}
\date{stacks-project:spaces-over-fields.tex:local-snapshot-20261012; GNU FDL 1.2 or later}
\begin{document}
\maketitle
\noindent\textit{Editorial note.} All delivered mathematical text is the author's own, copied from the pinned Stacks Project chapter file spaces-over-fields.tex; no statement or proof step was rewritten, reordered or removed.\par\medskip
\noindent\textit{Editorial note.} The counted proof is the author's complete argument: the reduction to a finite surjective chart with the dimension bound at the points over x, the construction of an invariant affine open containing the finite orbit, and the descent of that affine open to an affine open subspace through the submersiveness of the finite surjective map.\par\medskip
\noindent\textit{Editorial note.} The label lemma-codim-1-point-in-schematic-locus names the counted statement, and the proof environment immediately following it in the source is delivered verbatim as its proof.\par\medskip
\noindent\textit{Editorial note.} The counted statement is carried by the author's own lemma environment in the chapter source.\par\medskip
\noindent\textit{Editorial note.} Pre-open gate: python operations/agent-authoring/tools/precard\_gate.py --source spaces-over-fields.tex --label lemma-codim-1-point-in-schematic-locus --start 1659 --end 1701 --local-id 28060 returned RESULT=OK (L28060 free, occupancy FREE, public interval clear). Candidate taken from sources/stacks-spaces-over-fields/candidates-halfA.json.\par\medskip
\noindent\textit{Editorial note.} The statement opens with a native slogan environment stating the content informally; that environment is part of the delivered statement and is covered by the preamble span 8-13 of the chapter, which this card includes.\par\medskip
\noindent\textit{Editorial note.} The statement and proof are copied from the pinned chapter file; the cited inputs are kept as external pointers with their canonical Stacks tags.\par\medskip
\noindent\textit{Allowed standard prerequisites (editorial).}\begin{itemize}
\item Existence of a finite surjective scheme over the space. Limits of Spaces, the proposition producing a scheme finite and surjective over a given algebraic space. Supplies the finite surjective morphism U -> X used to build the groupoid presentation.
\item Dimension of the local rings of the cover. Decent Spaces, the lemma comparing the dimensions of the local rings of a point of the cover and of its image under a quasi-finite morphism. Supplies the bound on the dimensions of the local rings at the points over x.
\item Existence of an invariant affine open containing a finite orbit of dimension at most one. More on Groupoids, the lemma producing an invariant affine open containing a given orbit at points of dimension at most one. Supplies the invariant affine open W containing the points mapping to x.
\item The quotient of an invariant affine open is affine, and closed surjective maps are submersive. Cohomology of Spaces, the lemma identifying the image of an invariant affine open under a finite surjective morphism as an affine scheme, and Topology, the lemma that a closed surjective continuous map is submersive. Supply the open subspace X' which is an affine scheme and the openness of the image of W.
\end{itemize}\medskip
\ifdefined\hypertarget\hypertarget{sourcepointer-1}{}\fi
\noindent\textit{Source reference.} Original lemma 0ABW: sources/stacks-chow/tags.txt, source lines 13034--13034. Prerequisite pointer invoked at proof line(s) 1685; it names a standard fact of a sibling chapter that the counted argument uses, and every step of the counted proof built on it is retained verbatim in the delivered body.\par\medskip
\ifdefined\hypertarget\hypertarget{sourcepointer-2}{}\fi
\noindent\textit{Source reference.} Original lemma 0ABI: sources/stacks-chow/tags.txt, source lines 13021--13021. Prerequisite pointer invoked at proof line(s) 1689; it names a standard fact of a sibling chapter that the counted argument uses, and every step of the counted proof built on it is retained verbatim in the delivered body.\par\medskip
\ifdefined\hypertarget\hypertarget{sourcepointer-3}{}\fi
\noindent\textit{Source reference.} Original lemma 07VP: sources/stacks-chow/tags.txt, source lines 10014--10014. Prerequisite pointer invoked at proof line(s) 1699; it names a standard fact of a sibling chapter that the counted argument uses, and every step of the counted proof built on it is retained verbatim in the delivered body.\par\medskip
\ifdefined\hypertarget\hypertarget{sourcepointer-4}{}\fi
\noindent\textit{Source reference.} Original proposition 09YC: sources/stacks-chow/tags.txt, source lines 12560--12560. Prerequisite pointer invoked at proof line(s) 1678; it names a standard fact of a sibling chapter that the counted argument uses, and every step of the counted proof built on it is retained verbatim in the delivered body.\par\medskip
\ifdefined\hypertarget\hypertarget{sourcepointer-5}{}\fi
\noindent\textit{Source reference.} Original definition 04NA: sources/stacks-chow/tags.txt, source lines 6071--6071. Prerequisite pointer invoked at proof line(s) 1668; it names a standard fact of a sibling chapter that the counted argument uses, and every step of the counted proof built on it is retained verbatim in the delivered body.\par\medskip
\ifdefined\hypertarget\hypertarget{sourcepointer-6}{}\fi
\noindent\textit{Source reference.} Original lemma 0AAU: sources/stacks-chow/tags.txt, source lines 12997--12997. Prerequisite pointer invoked at proof line(s) 1694; it names a standard fact of a sibling chapter that the counted argument uses, and every step of the counted proof built on it is retained verbatim in the delivered body.\par\medskip
\begin{lemma}
\label{lemma-codim-1-point-in-schematic-locus}
\begin{slogan}
Separated algebraic spaces are schemes in codimension 1.
\end{slogan}
Let $S$ be a scheme. Let $X$ be an algebraic space over $S$.
Let $x \in |X|$. If $X$ is separated, locally Noetherian, and
the dimension of the local ring of $X$ at $x$ is $\leq 1$
(Properties of Spaces, Definition
\ref{spaces-properties-definition-dimension-local-ring}),
then there exists an open subspace of $X$ containing $x$ which is a scheme.
\end{lemma}

\begin{proof}
(Please see the remark below for a different approach avoiding the material on
finite groupoids.) We can replace $X$ by an quasi-compact neighbourhood of
$x$, hence we may assume $X$ is quasi-compact, separated, and Noetherian.
There exists a scheme $U$ and a finite surjective morphism $U \to X$,
see Limits of Spaces, Proposition
\ref{spaces-limits-proposition-there-is-a-scheme-finite-over}.
Let $R = U \times_X U$. Then $j : R \to U \times_S U$ is an equivalence
relation and we obtain a groupoid scheme $(U, R, s, t, c)$ over $S$
with $s, t$ finite and $U$ Noetherian and separated.
Let $\{u_1, \ldots, u_n\} \subset U$ be the set of points mapping to $x$. 
Then $\dim(\mathcal{O}_{U, u_i}) \leq 1$ by
Decent Spaces, Lemma
\ref{decent-spaces-lemma-dimension-local-ring-quasi-finite}.

\medskip\noindent
By More on Groupoids, Lemma
\ref{more-groupoids-lemma-find-affine-codimension-1}
there exists an $R$-invariant affine open $W \subset U$ containing
the orbit $\{u_1, \ldots, u_n\}$. Since $U \to X$ is finite surjective
the continuous map $|U| \to |X|$ is closed surjective, hence
submersive by Topology, Lemma
\ref{topology-lemma-closed-morphism-quotient-topology}.
Thus $f(W)$ is open and there is an open subspace $X' \subset X$
with $f : W \to X'$ a surjective finite morphism.
Then $X'$ is an affine scheme by
Cohomology of Spaces, Lemma
\ref{spaces-cohomology-lemma-image-affine-finite-morphism-affine-Noetherian}
and the proof is finished.
\end{proof}

\end{document}
